% !PRINTED 249
\weberpage{249}
Wenn man die Null mit zu den Quadraten zählt, so erhält man den Satz:

\gesperrt{6. Wenn $p = 2$ ist, so sind alle Zahlen in $\mathfrak{C}$ Quadrate. Ist $p$ ungerade, so giebt es $\frac{1}{2}(p^n + 1)$ Quadrate, $\frac{1}{2}(p^n - 1)$ Nichtquadrate in $\mathfrak{C}$.}

Bei ungeradem $p$ muss eine primitive Wurzel immer ein Nichtquadrat sein und in der Reihe (8) sind die Zahlen mit geraden Exponenten die Quadrate, die mit ungeraden Exponenten die Nichtquadrate.

Das Product und der Quotient zweier Quadrate ist offenbar wieder ein Quadrat.

Daraus folgt, dass das Product aus einem Nichtquadrat und einem von Null verschiedenen Quadrate ein Nichtquadrat ist. Denn ist das Product $\alpha\beta$ ein Quadrat und der eine Factor $\alpha$ ein Quadrat, so muss auch $\beta = \alpha\beta : \alpha$ ein Quadrat sein.

Weiter schliesst man daraus, dass das Product von zwei Nichtquadraten ein Quadrat ist. Denn bedeutet $\beta$ irgend ein Nichtquadrat, so lasse man in dem Producte $\beta\xi$ den Factor $\xi$ die sämmtlichen von Null verschiedenen Zahlen von $\mathfrak{C}$ durchlaufen. Das Product durchläuft dann dieselbe Zahlenreihe, und da für alle quadratischen $\xi$ das Product $\beta\xi$ Nichtquadrat ist, so muss es für die nichtquadratischen $\xi$ Quadrat sein.

Alles dies ergiebt sich auch sehr einfach aus der Darstellung (8) der Zahlen von $\mathfrak{C}$ durch eine primitive Wurzel.

Wir schliessen diese allgemeinen Betrachtungen mit dem Satze, den wir später brauchen werden:

\gesperrt{7. Jede nicht quadratische Zahl ist die Summe von zwei Quadraten.}

Hierbei ist $p$ als ungerade vorauszusetzen, da es nur dann nichtquadratische Zahlen giebt.

Es lässt sich dann zunächst jede Zahl $\beta$ als Differenz zweier Quadrate darstellen, wie die Identität

\[
\left(\beta + \tfrac{1}{4}\right)^2 - \left(\beta - \tfrac{1}{4}\right)^2 = \beta
\]

zeigt. Ist nun $-1$ ein Quadrat in $\mathfrak{C}$, so setze man

\[
\left(\beta + \tfrac{1}{4}\right)^2 = \xi^2, \quad -\left(\beta - \tfrac{1}{4}\right)^2 = \eta^2,
\]

und erhält

\[
\beta = \xi^2 + \eta^2.
\]

Ist aber $-1$ und also auch $p - 1$ Nichtquadrat, so suche man in der Reihe der natürlichen Zahlen $1, 2, \ldots, p - 1$, deren
%CONT>>
